% EM structure and teaching revision 20260928
\subsection{把拓扑、方程与功率检查分开}
设连通电路有 $N$ 个节点和 $B$ 条支路。固定一个参考电势后，可独立写出 $N-1$ 个节点电流方程；若改用支路电流未知量，则还需 $B-N+1$ 个独立回路方程及元件关系。逐个面孔写回路虽方便，但非平面画法中的线条交叉未必表示新节点。
例如节点电势 $V$ 经 $R_1$ 接到已知电势 $E$、经 $R_2$ 接地，同时有恒流源 $I_s$ 向该节点注入，则
\[
 \frac{V-E}{R_1}+\frac V{R_2}-I_s=0,
 \qquad V=\frac{E/R_1+I_s}{1/R_1+1/R_2}.
\]
令 $I_s=0$ 应退回分压公式；令 $R_2\to\infty$ 应得到 $V=E+I_sR_1$。这类极限检查比单纯把计算结果代回原方程更容易发现参考方向选错。
若理想电压源连接两个未知电势节点，把两节点作为超节点写总电流平衡，再补电压差约束。解完后分别计算各电阻消耗和各电源输出，检查全电路净功率为零；这一步同时检验数值、支路方向和正负号。

